\section{Statement of Compliance}

This trial will be conducted in accordance with the protocol, the principles of
the International Council for Harmonisation (ICH) Good Clinical Practice (GCP)
guideline E6(R3), and all applicable United States regulations governing both
the investigational drug and the investigational device, together with the
Physical AI overlay that governs the on-premises large language model (LLM)
directed robotic system. This Phase 2 study is the randomized, controlled,
multicenter efficacy evaluation that the first-in-human combined Investigational
New Drug (IND) and Investigational Device Exemption (IDE) protocol explicitly
deferred \cite{kawchak2026phase1}; the predicate study established the
daraxonrasib recommended Phase 2 dose (RP2D) and device feasibility and safety,
and this protocol now advances both components into a powered comparative trial.
Two distinct but harmonized regulatory spines apply concurrently across the eight
participating high-volume hepatopancreatobiliary (HPB) centers; both are carried
in full below, and neither is subordinated to the other. The Sponsor-Investigator
attests that the conduct of every study activity, the protection of every
enrolled participant, and the integrity of every recorded datum will satisfy the
more stringent of the two spines wherever they differ.

\subsection{Drug Arm Compliance (IND, 21 CFR Part 312)}

The perioperative daraxonrasib (RMC-6236) component is conducted under an IND in
accordance with Title 21 of the Code of Federal Regulations. Specifically, the
study complies with 21 CFR part 50 (Protection of Human Subjects, including the
informed-consent requirements of \S50.20 through \S50.27) \cite{ecfr2026part50};
21 CFR part 54 (Financial Disclosure by Clinical Investigators)
\cite{ecfr2026part54}; 21 CFR part 56 (Institutional Review Boards); and 21 CFR
part 312 (Investigational New Drug Application), including the sponsor and
investigator responsibilities of subparts B and D, the safety-reporting
requirements of \S312.32, and the IND-application and amendment provisions of
\S312.20 through \S312.31. Because the predicate Phase 1 study established the
RP2D \cite{kawchak2026phase1}, the drug arm is no longer a 3+3 dose-finding
evaluation; daraxonrasib proceeds at the fixed RP2D of 300 mg orally once daily
as a Phase 2 controlled efficacy evaluation under \S312.21(b). The IND remains in
effect following the review period of \S312.40, the protocol advance is filed as
an amendment under \S312.30; no participant is dosed before amended IND
is in effect.

\subsection{Device Arm Compliance (IDE, 21 CFR Part 812)}

The on-premises LLM-directed eight-arm robotic pancreaticoduodenectomy (Whipple)
system (PancreSpeed II) is an investigational medical device evaluated under an
IDE in accordance with 21 CFR part 812 \cite{ecfr2026part812}. The
Sponsor-Investigator and the reviewing Institutional Review Board (IRB) have
determined that the device presents a significant risk within the meaning of
\S812.3(m), because it is a surgical-intervention device whose patient-contact
motion, force application, vascular-exclusion gating, and emergency-stop behavior
present a potential for serious risk to the health, safety, or welfare of a
participant. The study therefore proceeds only under an IDE approved by both the
FDA and the IRB. The device arm additionally complies with 21 CFR part 50
\cite{ecfr2026part50} and 21 CFR part 56 for human-subject protection and IRB
oversight, and with 21 CFR part 11 (Electronic Records; Electronic Signatures)
\cite{ecfr2026part11} for the validation, audit-trail, and electronic-signature
controls that underpin the federated hash-chained command and telemetry record
described in \S10. Having met the early-feasibility objectives of the predicate
study \cite{kawchak2026phase1}, the device component is now framed as a
multicenter randomized controlled clinical study under part 812, in which the
comparative clinical performance of the device-directed operation is evaluated
against institutional-standard high-volume pancreaticoduodenectomy.

\clearpage

\subsection{Physical AI Overlay (21 CFR Part 312, Subpart J)}

The autonomy-bearing behavior of the LLM-directed system is governed by the
Physical AI overlay codified in 21 CFR part 312, Subpart J (\S312.400 through
\S312.405), as adapted for end-to-end Physical AI oncology investigations
\cite{kawchak2026adapt312}. For this Phase 2 study the Phase 0 gate is upgraded
relative to the predicate. As preconditions to randomized enrollment, Subpart J
now requires at least 5000 simulated procedures reproduced across at least three
independent simulation frameworks, with a tightened sim-to-real correspondence of
trajectory error $< 1.5$ mm and tip-force error $< 0.4$ N; a documented Unified
Safety Level (USL) readiness rating of $\geq 8.0$ for surgical deployment, raised
from the predicate threshold; per-site qualification together with multicenter
fleet harmonization so that every device instance behaves identically across the
eight centers; and a federated hash-chained audit trail spanning all sites under
21 CFR part 11 \cite{ecfr2026part11}. The system remains classified as a Class II
collaborative device under continuous human oversight with a guaranteed
emergency-stop latency of $\leq 500$ ms, operated as Stage 2 supervised
semiautonomy with immediate manual override; full autonomy is prohibited under 21
CFR \S312.21(e). These overlay requirements operate in addition to, and never in
place of, the part 312 and part 812 obligations above. The hash-chained audit
trail required under \S312.404 satisfies the part 11 electronic-records standard
and renders every LLM advisory, every robot command, and every safety-limit event
re-traceable to a deterministic seed and a specific repository commit.

\subsection{Good Clinical Practice and Electronic Records}

The trial adheres to ICH E6(R3) GCP, the international ethical and scientific
quality standard for designing, conducting, recording, and reporting clinical
investigations involving human participants. All electronic clinical data,
source records, and the cyber-physical telemetry stream are maintained under 21
CFR part 11 \cite{ecfr2026part11}, with validated systems, secure time-stamped
audit trails that do not obscure prior entries, role-based access controls, and
electronic signatures bound to the records they authenticate. The federated
hash-chained record specified by the Physical AI overlay provides cryptographic
continuity across the screening, operative, and follow-up phases and across all
eight participating sites, so that the multicenter telemetry corpus forms a single
tamper-evident chain of custody.

\subsection{Multicenter Single IRB and Co-Investment Governance}

Because this is a cooperative multicenter study, a single IRB of record (sIRB)
reviews the protocol for all eight participating sites in accordance with 45 CFR
46.114 \cite{commonrule_sirb}, with site-level context handled through reliance
agreements coordinated by the Coordinating Center. All co-investment in the trial
is disclosed and managed under 21 CFR part 54 \cite{ecfr2026part54}. Contributions
to the Patient-Aligned Co-Investment Facility (PACIF) are held behind a capital
firewall that isolates every funder from randomization, endpoint definition,
outcome adjudication, analysis, and publication; capital may purchase only
operational capacity and never scientific outcomes. This separation is governed
under the H.R. 9510 verification, validation, and uncertainty quantification
(VVUQ) financial-data standard \cite{kawchak2026hr9510}, which subjects the flow
and application of co-investment funds to the same auditable VVUQ discipline
applied to the trial's technical evidence.

\subsection{Trial Registration}

The study will be registered on ClinicalTrials.gov before enrollment of the
first participant. It is registered as an Interventional study with a randomized,
parallel-assignment allocation. The drug component is classified as Phase 2
(controlled efficacy evaluation of daraxonrasib at the established RP2D), while
the device component carries Phase: Not Applicable, consistent with the
ClinicalTrials.gov convention that FDA-defined phase categories apply to drugs and
biologics rather than to devices. The Primary Purpose is recorded as Treatment.
The study is open-label with blinded independent central review (BICR) of imaging,
central masked pathology, and blinded endpoint adjudication, and it carries
prespecified group-sequential interim analysis with efficacy and futility
stopping rules.

\subsection{Statement of Compliance Pathway}

Figure 1 renders the combined IND and IDE compliance spine with the Physical AI
Subpart J overlay as upgraded for Phase 2. The daraxonrasib drug arm (IND, part
312, now at the established RP2D) and the robotic Whipple device arm (IDE, part
812, significant-risk per \S812.3(m), now a randomized controlled study) converge
through the upgraded Phase 0 simulation-validation gate, the raised USL readiness
verification, the single IRB across the eight sites, and the Class II
collaborative classification into a single multicenter randomized enrollment
pathway.

\subsection{Signatures}

The Sponsor-Investigator and the Coordinating Principal Investigator have reviewed
and approved this protocol and confirm that the study will be conducted in
compliance with the protocol, ICH E6(R3) GCP, and the federal regulations cited
above. By signature below, each signatory agrees to conduct or supervise the study
in accordance with the protocol and all applicable requirements.

\vspace{0.6em}

\noindent\begin{tabularx}{\textwidth}{L{8.2cm} Y}
\toprule
\textbf{Sponsor-Investigator} & \textbf{Coordinating Principal Investigator} \\
\midrule
\vspace{0.1em}\rule{6.0cm}{0.4pt} & \vspace{0em}\rule{6.0cm}{0.4pt} \\
Kevin Kawchak & Name (printed) \\
CEO, ChemicalQDevice & Title / Institution \\[0.8em]
Signature and date: \rule{2.6cm}{0.4pt} & Signature and date: \rule{2.6cm}{0.4pt} \vspace{0.05cm} \\
\bottomrule
\end{tabularx}

\vspace{0.4em}

\noindent\begin{tabularx}{\textwidth}{L{8.2cm} Y}
\toprule
\textbf{Multi-Site Principal Investigator} & \textbf{Site / Institution} \\
\midrule
\vspace{0.1em}\rule{6.0cm}{0.4pt} & \vspace{0em}\rule{6.0cm}{0.4pt} \\
Name (printed) & One signature line per participating site \\
Title / Institution & (eight high-volume HPB centers) \\[0.8em]
Signature and date: \rule{2.6cm}{0.4pt} & Signature and date: \rule{2.6cm}{0.4pt} \vspace{0.05cm} \\
\bottomrule
\end{tabularx}

\vspace{0.3em}

\noindent The signatories acknowledge that this is an independent research draft;
it is not an active IND or IDE and is not endorsed by the FDA, NIH, HHS, an IRB,
ICH, or any sponsor.


\begin{mermaidfig}
\node[mmin]  (pre)   at (0,0)       {IND in effect\\daraxonrasib at RP2D\\Phase 2 efficacy};
\node[mmstep](ind)   at (0,-2.2)    {Amended IND\\\S312.21(b)\\300 mg PO daily};
\node[mmin]  (sr)    at (4.6,0.2)      {Significant-risk\\IDE \S812.3(m)\\randomized controlled};
\node[mmstep](ide)   at (4.6,-2.2)   {IDE approved\\(FDA plus IRB)\\PancreSpeed II};
\node[mmstep](p0)    at (2.3,-4.3)   {Phase 0 sim validation\\$\geq 3$ frameworks\\$<1.5$ mm / $<0.4$ N};
\node[mmstep](usl)   at (2.3,-6.3)   {USL readiness $\geq 8.0$\\fleet harmonization\\federated audit};
\node[mmgoal](cls)   at (2.3,-8.7)   {Single IRB (sIRB)\\8 HPB sites\\45 CFR 46.114};
\node[mmgoal](enr)   at (2.3,-10.9)  {Multicenter randomized\\enrollment 1:1\\Class II, $\leq 500$ ms e-stop};
\draw[mmedge] (pre) -- (ind);
\draw[mmedge] (sr)  -- (ide);
\draw[mmedge] (ind.south) |- (p0.west);
\draw[mmedge] (ide.south) |- (p0.east);
\draw[mmedge] (p0)  -- (usl);
\draw[mmedge] (usl) -- (cls);
\draw[mmedge] (cls) -- (enr);
\end{mermaidfig}
\vspace{-0.7cm}
\figcaption{Figure 1. Combined IND / IDE pathway with upgraded
Physical AI overlay for Phase 2. The daraxonrasib drug arm (IND, part 312, at the
established RP2D) and the robotic Whipple device arm (IDE, part 812,
significant-risk per \S812.3(m), randomized controlled) converge through the
upgraded Subpart J overlay of Phase 0 validation across at least three
frameworks, USL readiness $\geq 8.0$ with multicenter fleet harmonization, single
IRB review across eight sites, and Class II collaborative classification into a
single multicenter randomized enrollment pathway.}
